\section{Data Collection Methodology}
\label{sec:collection}

\subsection{Collection Fleet and Sensor Configuration}
\label{ssec:sensor_config}

Data collection was conducted using three instrumented vehicle platforms (designated
\textit{Alpha}, \textit{Beta}, and \textit{Gamma} fleets), each carrying an identical
sensor suite to ensure cross-fleet consistency. The sensor configuration was designed
to provide complete 360\textdegree{} perception coverage at both close and long range:

\begin{itemize}[leftmargin=2em, noitemsep]
  \item \textbf{Cameras}: 8 synchronized cameras at 10 FPS and $1920 \times 1080$~px
        resolution; 1 front-facing camera (120\textdegree{} FOV), 2 front-angled cameras
        (70\textdegree{} FOV), 1 rear camera (120\textdegree{} FOV), 2 rear-angled cameras
        (70\textdegree{} FOV), 1 front fisheye (190\textdegree{} FOV), 1 rear fisheye
        (190\textdegree{} FOV).
  \item \textbf{LiDAR}: 1 rooftop spinning LiDAR (64-beam, 20~Hz, 200~m range) and
        4 corner solid-state LiDARs (128-beam, 25~Hz, 100~m range each) for near-range
        coverage and redundancy.
  \item \textbf{Radar}: 5 mmWave radar units (2 front, 2 rear, 1 centre-front) providing
        velocity-resolved detections up to 250~m, operating under all weather conditions.
  \item \textbf{GNSS/IMU}: Dual-antenna GNSS with RTK correction ($\pm$2~cm position accuracy)
        and a 6-DoF IMU at 100~Hz for precise ego-motion estimation.
  \item \textbf{HD Maps}: Proprietary lane-level HD maps with semantic attributes
        (lane type, speed limit, road marking, traffic sign) at $\pm$10~cm accuracy.
\end{itemize}

All sensors are hardware-timestamped and cross-calibrated using a 3D calibration target
board at the start and end of each recording session. Intrinsic and extrinsic calibration
parameters are stored per session to account for thermal drift and minor sensor displacement.

\subsection{Collection Locations}
\label{ssec:collection_locations}

\lucid{} was collected across six cities in three continents, deliberately chosen to
maximise traffic culture diversity, road infrastructure variability, and climatic range:

\begin{itemize}[leftmargin=2em, noitemsep]
  \item \textbf{Shanghai, China} (4,312 scenarios): Dense urban intersections,
        mixed traffic with high cyclist and e-scooter density, and diverse pedestrian
        behaviour. Collected across Puxi (central) and Pudong (suburban) districts.
  \item \textbf{Beijing, China} (1,847 scenarios): Ring-road highway scenarios and
        Beijing's distinctive hutong narrow-road navigation. Four-season coverage
        including heavy winter snow.
  \item \textbf{Singapore} (1,623 scenarios): Multi-ethnic pedestrian behaviour,
        tropical rain and haze events, and left-hand driving infrastructure.
  \item \textbf{London, United Kingdom} (1,789 scenarios): Left-hand driving,
        roundabout-intensive road network, mixed cyclist traffic, and frequent
        drizzle and fog conditions.
  \item \textbf{San Francisco, USA} (1,634 scenarios): Steep gradients, cable-car
        interactions, and fog-bank maritime weather conditions.
  \item \textbf{Munich, Germany} (1,642 scenarios): Autobahn high-speed driving,
        tram intersections, and winter ice/snow scenarios.
\end{itemize}

\subsection{Collection Protocol}
\label{ssec:collection_protocol}

Collection sessions were planned using a stratified sampling strategy to ensure
balanced coverage of scenario categories, weather conditions, and time of day.
Each collection driver was paired with a safety monitor. Sessions were scheduled
across all weekday types (Monday--Sunday), time slots (06:00--22:00 for day/dusk;
22:00--05:00 for night), and weather windows (precipitation events were actively
targeted using local meteorological forecasts).

Raw sensor streams were recorded at full sensor rate and time-stamped using a
Precision Time Protocol (PTP) hardware synchronisation network. Each session
produced an average of 4.7~GB of raw data per minute. Over the 18-month campaign,
collections totalled approximately 6.2~TB of raw sensor data before preprocessing.

\subsection{Data Preprocessing Pipeline}
\label{ssec:preprocessing}

Raw data underwent a six-stage preprocessing pipeline:

\begin{enumerate}[leftmargin=2em, noitemsep]
  \item \textbf{Sensor alignment}: Camera images are undistorted and georectified;
        LiDAR point clouds are merged into a unified sensor frame using calibrated
        extrinsics; radar detections are transformed to the ego-vehicle coordinate system.
  \item \textbf{Scenario segmentation}: Raw continuous recordings are segmented into
        12--120~second clips using a combination of (a) GPS-based location triggers,
        (b) event detection (\eg{} hard-brake events from IMU), and (c) manual annotations
        by the collection team, yielding 14,218 candidate scenarios.
  \item \textbf{Quality filtering}: Scenarios with persistent sensor failure (camera blur
        score $> 0.6$, LiDAR dropout $> 15\%$, GNSS accuracy $> 1$~m) are discarded
        (1,371 scenarios removed). Scenarios with fewer than three annotatable causal
        events are also filtered (removed 0 after manual review, as this criterion was
        addressed during scenario targeting).
  \item \textbf{Anonymisation}: Faces of pedestrians and licence plates are automatically
        blurred using a dedicated detection model with a recall target of $> 99.5\%$,
        verified by manual spot-checks on 5\% of frames.
  \item \textbf{3D object tracking}: A multi-modal detector and tracker generates
        preliminary 3D bounding boxes and track IDs for all traffic participants,
        serving as pre-annotations to assist human annotators.
  \item \textbf{Keyframe selection}: One representative ``annotation keyframe'' per
        3-second window is selected for scene-level description and agent-level labelling.
\end{enumerate}

After quality filtering, the final collection contains 12,847 scenarios available
for annotation.

\subsection{Ethical Considerations}
\label{ssec:ethics}

Data collection was conducted in full compliance with local privacy laws in all
six jurisdictions, including GDPR (Europe), PIPL (China), and CCPA (California).
The collection protocol was reviewed and approved by the institutional ethics
board (\#IRB-2024-187). All vehicle occupants provided written informed consent.
Pedestrians and other road users are not individually identifiable in the released
dataset, as faces and licence plates are anonymised as described above.
